\section{Verification and validation assessment}

The assessment separated three questions: was the intended calculation implemented, was its numerical solution sufficiently settled and conservative, and did it reproduce physical separator behaviour? Configuration checks and simple analytical tests addressed the first question. Residual, mass-balance, inventory and sensitivity checks addressed the second. Independent measurements would be required for the third.

\subsection{Implementation and numerical controls}

Reading back solver settings, checking the mesh and reopening selected saved solutions verified that the intended configuration and results had been retained. Python and PyFluent supported setup and extraction. Solver transcripts supplied iteration numbers when interface counters were stale. Source fields already expressed as per-cell mass flow were summed, rather than multiplied by cell volume again.

The main full-geometry Mixture baseline used SIMPLE, PRESTO pressure, second-order momentum/turbulence and QUICK volume fraction. Later controlled cases used Coupled; the steady-pool model used compressive volume-fraction reconstruction and implicit body force. Each comparison retained its actual pseudo-time and relaxation settings. Global Time Step in these steady calculations was numerical pseudo-time, not elapsed physical time.

A frictionless single-liquid duct independently checked the outlet-resistance implementation. For \(K=0\) or 9 and driving pressure 200 or 800 Pa, the analytical velocity was

\[U=\sqrt{\frac{2\Delta p}{\rho(1+K)}}.\]

Here, \(\Delta p\) is inlet total pressure minus external static pressure, \(\rho\) is liquid density and \(K\) specifies added outlet resistance. The unit term accounts for outlet kinetic head. The corrected test used a modest nonzero starting velocity and 500 iterations per case. Numerical acceptance used the final 100 iterations; speed and finite-value guards covered all 500. This isolated a boundary-law implementation, not multiphase separator hydraulics. Separate hydrostatic, sealed-boundary, interface and pressure-datum controls examined startup disturbances under their own conditions.

\subsection{Mass accounting, film transfer and storage}

Boundary fluxes were positive into the domain. For an EWF-off liquid calculation, the steady remainder was

\[R_L=\sum_b\dot m_{L,b}+\dot m_{L,\mathrm{source}},\qquad M_L=\int_V\rho_L\alpha_L\,\mathrm dV.\]

Removal sources were negative and counted once. Solver reports that already included sources were distinguished from boundary-only reports. Liquid, vapour and solver-reported mixture balances were normalised by their corresponding feeds. Mean absolute error used the mean of absolute sample remainders, not the absolute signed mean. Where DPM or film transfer was active, its signed contribution belonged in the relevant subsystem account.

For physical transient flow, the storage-aware remainder was

\[R_{L,\mathrm{transient}}=\sum_b\dot m_{L,b}+\dot m_{L,\mathrm{source}}-\frac{\Delta M_L}{\Delta t},\]

with any active inter-model transfers included. Steady inventory change per iteration remained a numerical-development measure. It was not divided by an EWF timestep to invent a physical bulk storage rate.

The local film mass balance compared measured inventory change with integrated transfers over the film-time interval:

\[E_f=\Delta M_f-\int_{t_0}^{t_1}\left(\dot m_{\mathrm{to\ film,net}}-\dot m_{\mathrm{film,out}}-\dot m_{\mathrm{direct\ drain}}\right)\,\mathrm dt.\]

The net source was the film mass gained after any release already included in the solver report. Release was therefore subtracted separately only when it had not already been included. Error normalisation followed the declared integrated source or transfer basis. This accounting allowed a growing film to conserve mass; stationary film additionally required negligible storage.

For the combined bulk, film and particle system, internal transfers must cancel. A film-only balance could not establish conservation of the combined system when bulk flow was held fixed. A local film pass with bulk flow updating also left the other mass balances to be checked. Film removal, collector removal and particle boundary fates were kept distinct.

\subsection{Convergence, comparison windows and uncertainty}

Assessment combined residuals, conservation, storage and routing. Appendix B gives the numerical criteria and windows for the principal tests; these were study tolerances, not estimates of physical uncertainty. It also defines the pressure averaging surfaces and weights, which differed between model families, and the historical pressure-drift calculation. Reaching an iteration count or film age was a separate execution outcome. Some later EWF runs continued after numerical acceptance limits ceased to control continuation, so completion did not mean those limits had been met. Table 3 summarises the interpretation checks.

\begin{longtable}{@{}>{\raggedright\arraybackslash}p{0.29000000\dimexpr\linewidth-4\tabcolsep\relax}>{\raggedright\arraybackslash}p{0.34000000\dimexpr\linewidth-4\tabcolsep\relax}>{\raggedright\arraybackslash}p{0.37000000\dimexpr\linewidth-4\tabcolsep\relax}@{}}
\caption{Numerical checks and their interpretation boundaries.}\label{tab:table-3}\\
\toprule
\textbf{Check} & \textbf{Assessment} & \textbf{Interpretation limit} \\
\midrule
\endfirsthead
\toprule
\textbf{Check} & \textbf{Assessment} & \textbf{Interpretation limit} \\
\midrule
\endhead
Bulk solution & Active residuals, phase/source balances, inventory and wrong-outlet flow & A small scaled residual alone did not establish steady conservation \\
Film solution & Inner-film tolerance, Courant, thickness, source integration and storage & A local balance pass did not establish stationary or complete-system flow \\
Mesh comparison & Between-mesh changes versus within-run drift and spatial differences & Equal iteration count or film age did not ensure equal numerical maturity \\
Particle tracking & Fate-mass reconciliation and sensitivity to tracking budget & Accounting for incomplete mass did not determine its eventual destination \\
\bottomrule
\end{longtable}

Mesh comparison used cell counts, boundary identities and available quality metrics. Pressure, velocity and inventory histories tested iterative stability before a spatial-discretisation claim. Matched wall resolution was not established. For field comparisons, cell velocities were volume-averaged into \(20\times70\times20\) common bins. Each shared bin was weighted by the smaller volume assigned by the two meshes. The resulting root-mean-square (RMS) vector difference was normalised by reference RMS velocity. Cross-mesh coverage was 86--91\% of reference volume, compared with 100\% for same-mesh iteration comparisons. A coarser \(10\times35\times10\) projection retained the qualitative conclusion. These were comparisons over supported bins, not whole-domain discretisation-error estimates. No Grid Convergence Index was inferred from unstable sequences. Runtime comparisons measured cost for the stated cases, not time to equivalent accuracy.

\subsection{Validation basis}

The Purnanto study supplied the model reconstruction reference. The later enthalpy comparison used the project's July tetrahedral results from Fluent 2024 R2, while the reconstruction used 2025 R2. The exact July case binary and tracking cap were not recovered, so agreement tested reproduction of recorded estimates rather than exact replication. Matched measurements of carryover, brine discharge and pressure loss, with compatible definitions and uncertainty, were unavailable. Implementation checks and numerical sensitivity were therefore assessed separately from physical accuracy.
